% !TeX root = ../papers/nc-p3.tex
\section{頂点指数二の有限窓と四変数Clifford代数の一致}
\label{sec:index-two-central-clifford}
\begin{theorem}\label{thm:index-two-central-clifford}
四次中心切断を持つ有限窓のうち、頂点指数2の全ての添字保持同型類は、
四次元空間 $V$ に対する次の四空間の $\mathrm{PGL}(V)$ 軌道にちょうど一致する。
\[
 W\subset\operatorname{Sym}^2V^*,\qquad
 \text{$W$ は基点無し},\qquad
 \dim\big(\mathbb P(W)\cap\{\operatorname{rank}q\le2\}\big)=0 .
\]
各辺で同じ関係 $W^\perp\subset\operatorname{Sym}^2V$ を取るのが標準形である。
全点で最小周期は1、反対切断の階数は6である。中心の三次元多様体は $\mathbb P(W)$ 上の
行列式四次曲面で分岐する二重被覆で、中心偏極は超平面の二倍である。
交叉の非横断点や階数1を含む場合もこの判定と同型条件に含まれる。
この軌跡は既約で同型類の次元は9である。
\end{theorem}
\begin{proof}[随伴行列による文字の同定]
前定理の $s=2$ を使う。$U$ 上で $\det E_{i+1}/\det E_i=L$、$2L=H$ であり、
$E_{i+2}\simeq E_i\otimes L$ が成立する。実際この二束は行列式が一致し、
前定理の一般切断上の相互Euler数と二超平面による降下の証明をそのまま適用できる。
$E_0,E_1$ と $L$ の順次の捩りで列全体を表示する。
階数二では $A:E_0\to E_1$ に随伴行列を対応させる操作は束の線型同型
\[
 \mathcal Hom(E_0,E_1)\longrightarrow\mathcal Hom(E_1,E_0\otimes L)
\]
である。二つの順序の積はともに $\det A$ 倍の恒等になる。
従って全切断の四文字を両側で同じ $V$ としてよく、各二次の対称積の像は
行列式の偏極 $\operatorname{Sym}^2V\to H^0(U,L)$ に入る。

この右辺の次元は高々4である。$U$ 上の直線束を $X$ に反射的に延ばす。
一般超平面 $S$ は特異点を避けた滑らかなK3で、$L|_S$ は豊富、平方4であるから
$h^0(S,L)=4$ である。$H^0(X,L-H)=0$ は $2L=H$ と豊富性から従う。
実際非零切断の二乗は反射的延長で $-H$ の切断になる。従って制限は単射である。
一方どの辺でも全関係は六次元である。十次元の対称積を上の四次元以下へ送る核が
関係に入るため、像は四次元、核は六次元で関係全体に等しい。
二つの順序で偏極行列式は全く同じだから、全ての辺で同一の六対称関係となる。

$W$ をこの核の対称双対における直交補とする。通常の一周期代数の二次双対は
$\operatorname{Sym}V^*/(W)$ であり、有限なので $W$ は基点を持たない。
逆に $W$ が基点無しなら四二次式は四変数での正則列となり、この可換二次完全交叉は
KoszulかつFrobeniusでHilbert関数は $1,4,6,4,1$ である。
例えば剰余体のTate分解で一次の外積生成元と二次関係に対応するホモロジー次数2・内部次数2の
分割冪生成元を取れば線型性が分かる。従って逆の二次双対は標準Koszul左右AS正則である。
可換双対の四次traceは全ての文字に対称なので四テンソルの三切断の階数は同じ6である。
点対応の式は二共変元 $p,p'$ が $p p'\in W$ を満たすことに等しい。
二次式の二因子分解の射は交換による有限商で、像は階数2以下の軌跡である。
よって問題の有限性条件を得る（四空間は余次元3の階数軌跡と必ず交わる）。

$W^*$ の次数2の四中心文字を置き、二つの $V$ の偏極積がこれらの線型結合になる
通常のClifford代数を作る。四中心は偏極積で生成され、商代数への標準全射を得る
（すなわち二次関係だけの代数からこのClifford代数への全射）。
Cliffordの順序付き平方自由単項式は係数多項式環上の自由基底である。
従ってHilbert級数は両側で $(1-t)^{-4}$ となり全射は同型である。
各二次中心の平方を四パラメータに取れる。標準周期は可換二次双対のため通常の四周期である。
四次中心には二次中心の十積が独立に含まれるので、二つの中心型のうち指数2になる。
特にこの判定の逆方向でも四次中心列の条件は自動である。
\end{proof}

\begin{proof}[残余同型と中心]
$R=W^\perp$ に非零純関係はない。有限辺の一般結果からも、基点無し条件から
$x^2\in R$ が $W$ の基点を与えることからも従う。
$g,h:V\to V'$ が $(g\otimes h)R\subset\operatorname{Sym}^2V'$ を満たすとする。
$g^{-1}h=T$ と書くと、$R$ の全元は $(T\otimes1)r=(1\otimes T)r$ を満たす。
異なる固有値があれば一般固有空間の異なる組の成分は零となる。
$3+1$ 分解なら $R$ と三空間の対称二乗との交わりは少なくとも五次元で純二乗を含む。
$2+2$ なら対角二ブロックの対称空間全体が六次元なのでやはり純二乗が入る。
他の分割では六次元もない。一固有値で非自明な場合は冪零部の冪を取り、
平方零の $N\ne0$ に替える。階数1なら $N=u f$ の表示で $r f\in ku$ となり、
$R\subset\operatorname{Sym}^2(\ker f)+k(u\odot z)$（$f(z)=1$）。
やはり三空間の対称二乗との交わりは五次元以上である。
階数2なら像と核が同じ二平面 $P$ で、補空間と $P$ を $N$ で同定する。
下側対角ブロックは零、交差ブロックは対称となり解空間は六次元で
$\operatorname{Sym}^2P$ を含む。これで $T$ はスカラーに限られる。
従って連続する文字基底は射影的に同一で、添字保持同型は $W$ の一つの射影基底変更に一致する。

四つの二次中心を座標として普遍四行対称行列 $Q$ を取る。
偶Clifford代数の体積元は四文字の反対称化で、偶積と可換、平方は行列式である
（一般非退化対角基底で計算し、多項式恒等式で戻す）。
一般行列が退化なら非零の体積元の平方が零となり、中心切断の素性・隅の可除性に反する。
非退化な一般点では偶中心はスカラーと体積元で張られる。
四中心多項式環上でもこの二元が自由基底である。
実際平方自由経路の基底で体積元の最高外積係数は1であり、
一般中心元からその最高係数倍を引けばスカラーだけが残る。
四次のVeroneseに移しても一般中心は拡大しない。一つの二次中心 $y$ を可逆にすれば
全ての偶経路は適当な $y$ の冪によって四の倍数次にできるからである。
整数環上の中心の同定も同じ基底比較で従う。さらに各頂点の同じ四の倍数次の中心元が
折り畳み全体の中心に延びることを確認する。スカラー多項式は全頂点で同じにし、
体積元は一文字と反可換なので頂点で交代符号を付ければよい。
一頂点への制限は合成可能な非零元の積の性質で単射である。
従って中心は $z^2=\det Q$ の二重被覆の二次Veroneseであり、偏極は主張通りである。

最後に基点無しと有限交叉は $\operatorname{Gr}(4,10)$ の非空開条件である。
一般四空間は横断的に階数2の滑らかな軌跡と出会い、階数1を避ける。
二十因子点の配置はその座標平方評価が単射なので二次超曲面に含まれない
（定理\ref{thm:finite-square-chart}）。これらを一点ずつ固定する射影変換の固有空間は
二超平面以内に含まれるため、一般の射影安定群は有限である。
ゆえに既約な同型類の次元は $4(10-4)-15=9$ となる。
\end{proof}
